\chapter{Conclusi\'on}
\label{cap:conc}

Recopilando los resultados relevantes en la tesis, logramos determinar el
n\'umero policiaco de las gr\'aficas de $k$ fichas de cualquier estrella $S_n$.
Si bien este \'ultimo es el resultado principal de la primera secci\'on de
resultados, no podemos dejar a un lado la utilidad del algoritmo de
desobstrucci\'on, el cual nos facilit\'o la descripci\'on de una estrategia
ganadora para $C$ con exactamente $c(\t{k}{S_n})$ polic\'ias. As\'imismo,
destacamos tambi\'en el descubrimiento de una estrategia ganadora para $R$ si
$C$ tiene menos de $c(\t{k}{S_n})$ polic\'ias.

Por otro lado, una ligera modificaci\'on en la mencionada estrategia ganadora
para $C$ nos permiti\'o determinar una cota superior para el n\'umero policiaco
de la gr\'afica de $k$ fichas de cualquier \'arbol. Adem\'as, nuestro resultado
acerca de los retractos de gr\'aficas de $k$ fichas de un \'arbol evidenci\'o
que dicha cota superior es justa cuando los \'arboles tiene grado m\'aximo
suficientemente respecto al n\'umero de fichas.

Para el resto de gr\'aficas de fichas de \'arboles con grado m\'aximo bajo,
logramos mejor la cota inferior para estrellas subdivididas de grado bajo dada
por el lema de retractos en la gr\'afica de fichas de un \'arbol, esto siempre y
cuando el n\'umero de fichas sea menor al n\'umero de emparejamiento de la
estrella subdividida. M\'as a\'un, si la estrella subdividida tiene a lo m\'as
dos ramas de longitud par, esta cota cubre todos los posibles valores de $k$.
Particularmente para las triadas, se obtuvo una cota inferior para todas las
posibles $k$ sin importar la paridad de la longitud de sus ramas. Dicha cota
inferior tambi\'en induce una estrategia de escape para $R$ en funci\'on de
otros algoritmos, como uno que construye emparejamientos m\'aximos dada una
gr\'afica.

Finalmente, revisamos el caso particular de la gr\'afica de Petersen, en la que
logramos acotar el n\'umero policiaco de $\t{2}{P}$ al dar una estrategia
ganadora para $C$ con $4$ polic\'ias.

M\'as all\'a de las cotas dadas para el n\'umero policiaco en los distintos
casos de estudio, confiamos en que el dise\~no de estrategias nos ayude a
comprender los problemas que quedan sin soluci\'on en este trabajo. La
intuci\'on que nos dejan las estrategias desarrolladas nos llevan a pensar en
las siguientes dos conjeturas.

\begin{conjetura}
    Para toda $k\geq 2$, $0<n<2k-1$ y $m\geq 1$ suficientemente grande para que
    $k\leq \left\lfloor\frac{\abs{V_{S_{n,m}}}}{2}\right\rfloor$, se tiene que
    \[
        c\paren{\t{k}{S_{n,m}}} = \left\lceil \frac{\left \lfloor \frac{n-1}{2}
        \right \rfloor + k + 1}{2} \right \rceil.
    \]
\end{conjetura}

\begin{conjetura}
    Si $P$ denota la gr\'afica de Petersen, entonces $c(\t{2}{P})=4$.
\end{conjetura}

A manera de resumen, se presenta una tabla con las cotas encontradas.

% \begin{table}[h]
%     \begin{center}
%     \begin{tabularx}{\textwidth}{>{\centering\arraybackslash}X 
%         | >{\centering\arraybackslash}X 
%         | >{\centering\arraybackslash}X | >{\centering\arraybackslash}X |
%         >{\centering\arraybackslash}X | >{\centering\arraybackslash}X }
%         \hline
%         & $c(\t{k}{S_n})$ & $c(\t{k}{S(m_1,\dots,m_n)})$ & $n=2$ & $T$
%         \\
%         \hline
%         $k< k_n$ & $k$ & $k$ & $k$ & $k$ \\
%         \hline
%         $k_n < k \leq \alpha'$ & N/A & $\left \rceil \frac{k+1}{2} \right\rceil$ & & \\
%         \hline
%         $\alpha' < k \leq \left\lfloor \frac{\abs{V}}{2} \right \rfloor$ & $k$ &
%         $\leq k$ & & 
%     \end{tabularx}
% \end{center}
% \end{table}